\begin{answer}
\begin{enumerate}[label=\roman*.]

\item \textbf{Using a different constant learning rate.} \textbf{No.} The root cause is that the optimal $\theta$ has infinite norm, not the magnitude of the learning rate. With any constant step size, the gradient never vanishes at a finite point for linearly separable data, so the algorithm still diverges.

\item \textbf{Decreasing the learning rate by $1/t^2$.} \textbf{Yes.} The total displacement of $\theta$ is bounded by $\sum_{t=1}^{\infty} \alpha_t \|\nabla\ell\| \leq C \sum_{t=1}^{\infty} \frac{1}{t^2} = \frac{C\pi^2}{6} < \infty$. Therefore $\theta$ is confined to a compact set and the step sizes $\alpha_t \|\nabla\ell_t\| \to 0$, so the stopping criterion $\|\theta^{(t+1)} - \theta^{(t)}\| < \varepsilon$ will eventually be satisfied.

\item \textbf{Linear scaling of the input features.} \textbf{No.} Linear rescaling is an invertible linear transformation; it does not change whether the data is linearly separable. The same divergence issue remains.

\item \textbf{Adding an $L_2$ regularization term $\|\theta\|_2^2$.} \textbf{Yes.} The penalized objective $\ell(\theta) - \lambda\|\theta\|_2^2$ is strongly concave and its maximizer satisfies $\nabla\ell(\theta^*) = 2\lambda\theta^*$. For linearly separable data the regularizer prevents $\|\theta\|\to\infty$ because the penalty grows without bound. The unique finite maximizer exists and gradient ascent converges to it.

\item \textbf{Adding zero-mean Gaussian noise to the training data or labels.} \textbf{Yes (with high probability).} With any continuous noise distribution, the probability that the perturbed dataset remains linearly separable is zero. Once the data is no longer separable, the log-likelihood has a finite maximum and the algorithm converges.

\end{enumerate}
\end{answer}
